%!TEX root = ../thesis.tex
\ifpdf
    \graphicspath{{Sections/5_discussion_figs/Raster/}{Sections/5_discussion_figs/PDF/}{Sections/5_discussion_figs/}}
\else
    \graphicspath{{Sections/5_discussion_figs/Vector/}{Sections/5_discussion_figs/}}
\fi

\chapter{Discussion}
\label{sec:discussion}

\section{Comparison with Hu et al.}

We retained the training and patch-score construction methods of Hu et al. \cite{hu_learning_2024} while moving from AAPM parallel-beam to LIDC-IDRI fan-beam data, allowing us to test the transfer of these methods across datasets and operators.

\subsection{Findings that Transferred}

Hu et al. reported PaDIS-VE-DPS PSNRs of $33.57\,\mathrm{dB}$ at 20 views and $29.48\,\mathrm{dB}$ at 8 views \cite{hu_learning_2024}. Our Paper-Faithful track reached $32.71$ and $27.15\,\mathrm{dB}$. Despite changed data and geometry, the 20-view result is within $0.86\,\mathrm{dB}$, and the 8-view result substantially improves on FDK and CP-TV. PaDIS also generated coherent whole CT slices from patches, whereas fixed stitching and averaging produced less credible anatomies, supporting use of randomly shifted patches.

At $512\times512$, LION-Physical PaDIS-VE-DPS obtained $40.08\,\mathrm{dB}$ PSNR and $0.960$ SSIM, significantly outperforming CP-TV and FDK, and supporting the high-resolution motivation \cite{hu_learning_2024}. PaDIS therefore remains useful where whole-image memory is prohibitive.

\subsection{Findings that Did Not Transfer}

We did not reproduce the claim that PaDIS consistently outperforms whole-image diffusion \cite{hu_learning_2024}. In the Paper-Faithful track, whole-image VE-DPS exceeded PaDIS-VE-DPS by $1.26\,\mathrm{dB}$ at 20 views and $1.95\,\mathrm{dB}$ at 8 views. The LION-Physical comparison showed the same reversal, also visible in \autoref{fig:results-ct20-additional} and \autoref{fig:results-ct8-additional}.

The whole-image prior was not, however, uniformly stronger. Validation-selected LION-Physical PaDIS-VE-DDNM achieved the best 20- and 60-view PSNRs and exceeded its whole-image counterpart by $2.30\,\mathrm{dB}$ at 20 views. Langevin and VE-DPS favoured the whole-image prior, while predictor-corrector rankings depended on the metric. These comparisons do not isolate architecture because settings were validation-selected per prior and sampler, but show that PaDIS performance depends on how its patch score is used.

\subsection{Ablation Findings}

We repeated Hu et al.'s patch-size, training-set-size and positional-encoding ablations \cite{hu_learning_2024} on LIDC-IDRI. The initialisation and schedule comparisons instead address paper--code differences.

Smaller patches gave no consistent benefit. The $96\times96$ model nearly matched whole-image VE-DPS, while the default $56\times56$ model was $1.43\,\mathrm{dB}$ lower. Non-monotonic improvement at $8\times8$ suggests that receptive field, optimisation, patch count and partition geometry all contribute. \autoref{tab:patch-size} and \autoref{fig:results-patch-size} favour global context when memory permits, although non-default models inherited the $P=56$ inference setting and may not be optimal. PaDIS's clearest advantage is therefore feasibility rather than quality from very small patches. Using all processed LIDC-IDRI slices was also not seen to improve PaDIS-VE-DPS under fixed within-prior training time.

Removing position channels reduced PSNR by only $0.3$ to $0.4\,\mathrm{dB}$, in contrast to the failure reported by Hu et al. \cite{hu_learning_2024}. The CT measurements and repeated data-consistency updates appear to provide enough global information to locate local anatomy. This finding remains conditional on the shared inference setting and does not establish that positional channels are unnecessary for unconditional generation, where spatial context cannot be established by measurement-consistency. Similarly, FDK and Gaussian initialisation converged to almost identical results after 1000 model evaluations. The geometric noise schedule had a larger effect, improving PSNR by about $1.3\,\mathrm{dB}$ relative to the original repo EDM-style default.

\subsection{Original PaDIS Specification and Implementation}
\label{sec:discussion-reproducibility}

We treated the released code as part of the specification of the original PaDIS work rather than as a separate implementation \cite{hu_learning_2024,hu_jasonhu4padis_2026-1}. Together, the paper and code leave choices ambiguous. The variable-size U-Net is attributed to different references, while the number of sampled patches, validation split, checkpoint rule and selection of 25 test slices are not reported. View and detector counts are given without the physical field of view, detector spacing or complete angular conventions needed to define $A$. The released 180-view fan-beam configuration also spans only $180^\circ$, which is not explicit in the main description. Fine-tuning is reported without the search or selection protocol. Inconsistencies include their Table 1 not always bolding the best SSIM and their Table 7 describing a dataset-size analysis despite comparing samplers.

The implementations in the corresponding released repository resolve some details but differ from the algorithms in the original PaDIS paper \cite{hu_learning_2024,hu_jasonhu4padis_2026-1}. The main VE-DPS path starts from filtered back-projection rather than the Gaussian state in Algorithm 1 of the paper, differentiates the residual norm rather than its square and defaults to an EDM-style schedule. The predictor--corrector path evaluates the corrector denoiser at the current level, $\sigma_k$, but divides the denoising residual by $\sigma_{k+1}^2$ and reuses the patch layout drawn for the predictor. The paper reports 1000 neural function evaluations (NFEs) for all diffusion experiments, whereas the released CT command uses 100 noise levels. Their VE-DPS, Langevin, and DDNM implementations then retain 1000 evaluations through 10 updates per level. Their Predictor--corrector implementation instead makes two evaluations between adjacent levels, totalling 198. The repository also advertises a CT checkpoint, but the Google Drive link returns a file does not exist error. Several hidden APIs also had to be re-integrated before the comparison methods in the repository could be run. As demonstrated in \autoref{tab:main-quality} and \autoref{tab:position}, these paper--code differences do materially change reconstruction.



\section{Extensions to the Original Work}

Paper-Faithful and Repo-Faithful represent two readings of the combined specification. LION-Physical is our extension and replaces geometry-specific calibration with an operator-derived scale. Its results assess transfer and stabilisation rather than direct reproduction of Hu et al.'s algorithms.

\subsection{Operator-Normalised Data Consistency}

The LION-Physical VE-DPS rows reached relative sinogram residuals near $10^{-3}$, over an order of magnitude below many Paper-Faithful and Repo-Faithful results. Diffusion models can otherwise produce plausible images that are not well supported by the measurements. In the limited-angle experiment, whole-image VE-DPS gave the best PSNR and SSIM, while PaDIS-VE-DPS gave the smallest residual. There is therefore a trade-off between similarity to the held-out target and strict agreement with the measured data.

The Repo-Faithful fan-beam updates required calibrated gradient scales of $0.0405$ and $0.1022$. Replacing them with operator normalisation based on $\lVert F\rVert_2^2$ still gave high-quality results across the angular settings. These constants therefore correct for a particular implementation rather than defining PaDIS itself. Operator normalisation is more reproducible because it remains meaningful when the geometry changes.

\subsection{Fixed Patch Assembly}

Fixed averaging obtained the highest reported SSIM and lowest MAE in parts of \autoref{tab:main-quality}, but these rows contain only four images and have wide uncertainty. Its unconditional samples also lacked credible global anatomy. These results show that overlap handling can suppress seams but do not establish an improvement over random PaDIS partitions.

\section{Limitations and Future Work}

As in the original CT experiments of Hu et al. \cite{hu_learning_2024}, our sinograms omit photon noise, scatter, beam hardening and motion. We therefore investigate angular undersampling rather than realistic low-dose CT. The experiments are two-dimensional, use one dataset and one training run per model, and evaluate 25 slices rather than independent scans. The fixed-overlap and native-resolution rows use four images. Bootstrap intervals measure variation across slices rather than retraining uncertainty. Some high-view methods inherit 20-view tuning, and we did not tune the patch-size and no-position rows per model. PSNR, SSIM and MAE also cannot determine whether a nodule was preserved or hallucinated.

Future work should add calibrated photon noise and task-level evaluation of clinically relevant features. Equal-epoch, multi-seed training would separate dataset size from optimisation time. Averaging scores over several partitions, or a lower-variance deterministic estimator, may help with accelerated samplers. External datasets and expert assessment are required before making clinical claims.

Overall, we reproduce the ability of PaDIS to assemble coherent, high-quality reconstructions from patch scores, including at native resolution. We do not reproduce universal superiority over whole-image diffusion. Instead, our results show where PaDIS is useful, which implementation choices change its ranking and how operator-derived scaling improves transfer between CT geometries.
